% ============================================================================
%  SECTION 7 — Stability capstone
% ============================================================================
\begin{frame}[t]{}
  \vfill\begin{center}{\Huge\color{YaleBlue} 7.\ \ Stability}\\[0.5em]
  {\color{BoxGray} adaptation vs.\ aneurysm --- the one big question}\end{center}\vfill
\end{frame}

\begin{frame}[t]{Where instability comes from: the Laplace exponent}
  Growth adds mass, \emph{lowering} required stress (stabilising); dilation
  \emph{raises} it, \textbf{quadratically} for the artery. Lose elastin and dilation
  wins the race --- the feedback turns positive: \textbf{runaway}.
  \slboxeq{2.2em}{eq:stab}{\left.\frac{\dd\bar\sigma}{\dd\theta}\right|_{\theta^*}<0
    \ \Leftrightarrow\ \text{tissue adapts}}
  \bigfigw{stab_criterion_map.png}{0.60\linewidth}{Stable / unstable map over insult
  strength --- the bar (linear) can never lose this race.}
  \figcredit{Fig.: A. Gebauer (TUM lecture); Cyron, Wilson, Humphrey, J R Soc Interface (2014)}
\end{frame}

\begin{frame}[t]{Same tissue, two insults}
  Read stability two ways: \textbf{existence} --- the equilibrated equation
  \eqref{eq:equil} has a root iff the tissue can adapt (no time integration); or
  \textbf{transient} --- integrate and watch it plateau or climb.
  \bigfigw{fig06_stability.pdf}{0.80\linewidth}{Retaining 30\% elastin the artery
  adapts (blue); retaining 3\% it dilates without bound (red). Map computed from the
  equilibrated existence test.}
  \figcredit{Reproduced with the \texttt{gr} package (this repo)}
\end{frame}

\begin{frame}[t]{How the four theories answer --- the takeaways}
  \small
  \begin{center}
  \renewcommand{\arraystretch}{1.4}
  \begin{tabular}{@{}p{0.24\textwidth}p{0.20\textwidth}p{0.48\textwidth}@{}}
    \toprule
    \textbf{Theory} & \textbf{Predicts stability?} & \textbf{What it says} \\
    \midrule
    \textbf{\color{YaleBlue}Kinematic growth} & No --- \emph{imposed} &
      Reaches the \emph{prescribed} homeostatic set-point, or runs away. Stability
      is built in, not learned from turnover. \\
    \textbf{\color{YaleBlue}Constrained mixture (full)} & Yes &
      Adapts or dilates \textbf{depending on the insult} --- the faithful
      reference. \\
    \textbf{\color{YaleBlue}Homogenized CMM} & Yes &
      Same verdict as the full model, following its \textbf{time trajectory}
      cheaply. \\
    \textbf{\color{YaleBlue}Equilibrated CMM} & Yes --- \emph{instantly} &
      \textbf{Matches} the transient theories \emph{if} a root exists; \emph{no}
      root $\Rightarrow$ unbounded growth. \\
    \bottomrule
  \end{tabular}
  \end{center}
  \vfill
  \begin{center}\color{AccentRed}
  Stability \emph{depends on the insult}, and --- with turnover-based theories ---
  it can be \emph{predicted}, not just assumed.
  \end{center}
\end{frame}

